\documentclass{article}
\usepackage{geometry}
\geometry{margin=1.5cm, vmargin={0pt,1cm}}
\setlength{\topmargin}{-1cm}
\setlength{\headheight}{12.5pt}
\setlength{\paperheight}{29.7cm}
\setlength{\textheight}{25.3cm}

% useful packages.
% \usepackage[UTF8]{ctex}
\usepackage{fancyhdr}
\usepackage{layout}
\usepackage{tikz}
\usetikzlibrary{arrows.meta}

% import common macros
% % % % % % % % % % % % % % % % % % % % % % % % % % % % % % % %
% Common Macros for LaTeX
% % % % % % % % % % % % % % % % % % % % % % % % % % % % % % % %

% Useful packages
\usepackage{amsfonts}
\usepackage{amsmath}
\usepackage{amssymb}
\usepackage{amsthm}
\usepackage{enumerate}
\usepackage{graphicx}
\usepackage{multicol}
\usepackage[ruled,linesnumbered]{algorithm2e}

% Math operators and common symbols
\newcommand{\dif}{\mathrm{d}}
\newcommand{\innerProd}[1]{\left\langle #1 \right\rangle}
\newcommand{\difFrac}[2]{\frac{\dif #1}{\dif #2}}
\newcommand{\pdfFrac}[2]{\frac{\partial #1}{\partial #2}}
\newcommand{\T}{\mathrm{T}}
\newcommand{\norm}[1]{\left\| #1 \right\|}
\newcommand{\abs}[1]{\left| #1 \right|}

% Vector and Matrix shortcuts
\newcommand{\vecTwo}[2]{
  \begin{bmatrix}
    #1 \\
    #2
\end{bmatrix}}
\newcommand{\vecThree}[3]{
  \begin{bmatrix}
    #1 \\
    #2 \\
    #3
\end{bmatrix}}
\newcommand{\vecTwoH}[2]{
  \begin{bmatrix}
    #1 & #2
  \end{bmatrix}
}
\newcommand{\vecThreeH}[3]{
  \begin{bmatrix}
    #1 & #2 & #3
  \end{bmatrix}
}

% Common fields
\newcommand{\R}{\mathbb{R}}
\newcommand{\C}{\mathbb{C}}
\newcommand{\Z}{\mathbb{Z}}
\newcommand{\N}{\mathbb{N}}

% Solutions and proofs
\newenvironment{solution}{
\begin{proof}[解]}{
\end{proof}}

% Utility to get environment variables (requires lualatex)
\newcommand{\getEnv}[2]{\directlua{tex.print(os.getenv("#1") or "#2")}}


% Name and uID
\newcommand{\name}{\getEnv{NAME}{NAME}}
\newcommand{\uid}{\getEnv{UID}{UID}}
\newcommand{\course}{Complex Analysis}
\newcommand{\hwNum}{5}

\begin{document}

\pagestyle{fancy}
\fancyhead{}
\lhead{\zhstr{\name} (\uid)}
\chead{\course\ \#\hwNum}
\rhead{\today}

\section{Exercise 3.39}

\textbf{Theorem 2.73 (original rectangular-domain version).}
Suppose $u:D\to\R$ is harmonic, where
$D:=(a,b)\times(c,d)\subset\C$ is an open rectangle with sides
parallel to the coordinate axes. Then there exists an analytic
function $f:D\to\C$ with $u$ as its real part, and its harmonic
conjugate is uniquely determined up to an additive constant by
\[
\begin{aligned}
v(x,y)
&:=\int_{y_0}^{y}\partial_xu(x,t)\,dt
-\int_{x_0}^{x}\partial_yu(t,y_0)\,dt
+v(x_0,y_0),
\end{aligned}
\]
where $x_0\in(a,b)$ and $y_0\in(c,d)$.

Prove that Theorem 2.73 holds not only for rectangular
domains but also for simply connected domains.

\begin{solution}
	Let $u\in C^2(D)$ be harmonic on a simply connected domain
	$D\subset\C$, and define
	\[
		g(z):=u_x(x,y)-i\,u_y(x,y),\qquad z=x+iy.
	\]
	Then $g\in C^1(D)$ and
	\[
		\partial_x(\Re g)=u_{xx}=-u_{yy}=\partial_y(\Im g),
		\qquad
		\partial_y(\Re g)=u_{xy}=-\partial_x(\Im g),
	\]
	so $g$ satisfies the Cauchy--Riemann equations; hence $g$ is
	analytic on $D$.

	Because $D$ is simply connected, $g$ has a primitive $G$ on $D$:
	$G'=g$.  Write $G=U+iV$.  For analytic $G$, we have
	$G'=U_x-iU_y$, therefore
	\[
		U_x=u_x,\qquad U_y=u_y.
	\]
	So $\nabla(U-u)=0$, and on the connected domain $D$ this implies
	$U-u\equiv C$ for some real constant $C$.
	Set
	\[
		f:=G-C=(U-C)+iV=u+iV.
	\]
	Then $f$ is analytic on $D$ and has real part $u$, so $V$ is a
	harmonic conjugate of $u$.

	Uniqueness up to an additive constant: if
	$f_1=u+iv_1$ and $f_2=u+iv_2$ are analytic, then
	$h:=f_1-f_2=i(v_1-v_2)$ is analytic with zero real part.
	By Cauchy--Riemann, $\nabla(\Im h)=0$, hence $v_1-v_2$ is constant.
	Thus Theorem 2.73 extends from rectangles to simply connected
	domains.\qedhere
\end{solution}

\section{Exercise 3.41}

Determine whether the following functions have primitives on
the indicated domains and fully justify your answer.

(a) $f(z)=1/z$ on $\C^\bullet$;

(b) $g(z)=\bar{z}$ on $\C$.

\begin{solution}
	\textbf{(a)}\;
	If $F'=1/z$ on $\C^\bullet$, then by Theorem~3.16
	every closed-contour integral of $1/z$ in $\C^\bullet$ vanishes.
	But for $\gamma(t)=e^{it}$, $t\in[0,2\pi]$:
	\[
		\oint_\gamma\frac{dz}{z}
		=i\int_0^{2\pi}dt=2\pi i\neq0.
	\]
	Contradiction.

	\medskip

	\textbf{(b)}\;
	For the same $\gamma$:
	\[
		\oint_\gamma\bar z\,dz
		=\int_0^{2\pi}e^{-it}\cdot ie^{it}\,dt=2\pi i\neq0.
	\]
	By Theorem~3.16, $\bar z$ has no primitive on $\C$.\qedhere
\end{solution}

\section{Exercise 3.42}

Let $\gamma : [0,1] \to H := \{z \in \C : \Im z > 0\}$ be a
piecewise smooth curve with $\gamma(0)=1+i$ and $\gamma(1)=i$.

(a) Do $z$ and $1/z$ have primitives on $H$? Justify your answer.

(b) Calculate $\int_\gamma z\,dz$ and $\int_\gamma z^{-1}\,dz$.

\begin{solution}
	\textbf{(a)}\;
	$F(z)=z^2/2$ satisfies $F'=z$ on $H$.
	Since $0\notin H$ and $H$ is simply connected, the branch
	$\Log_H z:=\ln|z|+i\Arg_H z$ ($\Arg_H z\in(0,\pi)$)
	is single-valued on $H$ with $(\Log_H z)'=1/z$
	(cf.\ Corollary~3.15).

	\medskip

	\textbf{(b)}\;
	Both primitives exist on $H$; by Theorem~3.40 the integrals
	are path-independent:
	\begin{align*}
		\int_\gamma z\,dz
		&=\frac{z^2}{2}\bigg|_{1+i}^{i}
		=\frac{i^2}{2}-\frac{(1+i)^2}{2}
		=-\frac12-i,\\[6pt]
		\int_\gamma z^{-1}\,dz
		&=\Log_H z\bigg|_{1+i}^{i}
		=\Bigl(0+i\frac\pi2\Bigr)
		-\Bigl(\tfrac12\ln2+i\frac\pi4\Bigr)
		=-\frac12\ln2+\frac{\pi i}{4}.\qedhere
	\end{align*}
\end{solution}

\section{Exercise 3.51}

Prove that if $\{E_n\}$ is a sequence of closed nonempty and
bounded sets in a complete metric space $X$, if
$E_{n+1} \subset E_n$, and if
\[
	\lim_{n\to\infty} \operatorname{diam}(E_n)=0,
\]
then
\[
	\bigcap_{n=1}^{\infty} E_n
\]
consists of exactly one point.

\begin{solution}
	\textbf{Existence.}\;
	Pick $x_n\in E_n$.  For $m\ge n\ge N$,
	\[
		d(x_m,x_n)\le\operatorname{diam}(E_n)
		\le\operatorname{diam}(E_N)\to0,
	\]
	so $(x_n)$ is Cauchy; by completeness $x_n\to x$.
	For each $k$, $\{x_n:n\ge k\}\subset E_k$ and $E_k$ is closed,
	so $x\in E_k$.  Hence $x\in\bigcap_{n\ge1}E_n$.

	\medskip

	\textbf{Uniqueness.}\;
	If $x,y\in\bigcap_{n\ge1}E_n$, then
	$0\le d(x,y)\le\operatorname{diam}(E_n)\to0$,
	so $x=y$.\qedhere
\end{solution}

\section{Exercise 3.52}

Use the conclusion in Exercise 3.51 to prove the Baire
category theorem D.80.

\begin{solution}
	We prove: \emph{in a complete metric space $X$, every countable
	intersection of open dense sets is dense.}

	Let $\{U_n\}_{n\ge1}$ be open dense subsets.
	Fix any nonempty open ball $B_0:=B(x_0,r_0)$;
	we show $B_0\cap\bigcap_n U_n\neq\emptyset$.

	Construct closed balls $E_n:=\overline{B(x_n,r_n)}$
	inductively:
	\begin{itemize}\setlength\itemsep{2pt}
		\item \emph{Base} ($n=1$):\;
		$B_0\cap U_1$ is nonempty ($U_1$ dense) and open.
		Choose $E_1\subset B_0\cap U_1$ with $r_1\le2^{-1}$.
		\item \emph{Step} ($n\to n{+}1$):\;
		$\operatorname{int}(E_n)\cap U_{n+1}$ is nonempty
		($U_{n+1}$ dense) and open.
		Choose $E_{n+1}\subset E_n\cap U_{n+1}$
		with $r_{n+1}\le2^{-(n+1)}$.
	\end{itemize}
	Then $\{E_n\}$ are nonempty, closed, bounded, nested,
	and $\operatorname{diam}(E_n)\le2^{1-n}\to0$.
	By Exercise~3.51, $\bigcap_{n\ge1}E_n=\{x_*\}$.

	Since $E_n\subset E_1\subset B_0$ and $E_n\subset U_n$
	for all $n$:
	\[
		x_*\in B_0\cap\bigcap_{n\ge1}U_n,
	\]
	so $B_0\cap\bigcap_n U_n\neq\emptyset$.
\end{solution}

\section{Exercise 3.58}

For an analytic function $f : \C \to \C$ satisfying
\[
	\forall z\in\C,\quad f(z+1)=f(z),
\]
prove that the value of
\[
	\int_{[w,w+1]} f(\zeta)\,d\zeta
\]
is independent of the choice of $w\in\C$.

\begin{solution}
	Set $I(w):=\int_{[w,\,w+1]}f(\zeta)\,d\zeta$.
	Fix $w_0,w_1\in\C$ and let $\Gamma$ be the closed polygonal
	contour
	\[
		w_0\to w_1\to w_1{+}1\to w_0{+}1\to w_0.
	\]
	Since $f$ is entire, $\oint_\Gamma f\,dz=0$:
	\[
		\underbrace{\int_{w_0}^{w_1}f\,dz}_{A}
		+I(w_1)
		+\underbrace{\int_{w_1+1}^{w_0+1}f\,dz}_{B}
		-I(w_0)=0.
	\]
	By periodicity $f(\eta+1)=f(\eta)$ and substitution
	$\zeta=\eta+1$:
	\[
		\int_{w_0+1}^{w_1+1}f(\zeta)\,d\zeta
		=\int_{w_0}^{w_1}f(\eta+1)\,d\eta
		=\int_{w_0}^{w_1}f(\eta)\,d\eta=A,
	\]
	so $B=-\int_{w_0+1}^{w_1+1}f\,d\zeta=-A$.
	Hence $0=A+I(w_1)-A-I(w_0)$, i.e.\ $I(w_1)=I(w_0)$.\qedhere
\end{solution}

\section{Exercise 3.59}

For bounded simply connected domain $D$,

(a) prove
\[
	\oint_{\partial D} \overline{z}\,dz = 2i\,S(D),
\]
where $S(D)$ is the area of $D$;

(b) show that $f$ being analytic on $\bar{D}$ implies
\[
	\max_{z\in\partial D} |\bar{z}-f(z)| \ge 2\,\frac{S(D)}{\ell(\partial D)},
\]
where $\partial D$ is the boundary of $D$ and
$\ell(\partial D)$ its arc length.

\begin{solution}
	\textbf{(a)}\;
	$z=x+iy$ gives
	\[
		\bar z\,dz
		=(x-iy)(dx+i\,dy)
		=\underbrace{(x\,dx+y\,dy)}_{=\,d(x^2+y^2)/2}
		+i(x\,dy-y\,dx).
	\]
	The first term is exact, so its closed integral vanishes.
	Green's theorem on the second:
	\[
		\oint_{\partial D}(x\,dy-y\,dx)
		=\iint_D(1+1)\,dA=2\,S(D).
	\]
	Therefore $\oint_{\partial D}\bar z\,dz=2i\,S(D)$.

	\medskip

	\textbf{(b)}\;
	Set $h(z):=\bar z-f(z)$, \;$M:=\max_{\partial D}|h|$.
	By the ML inequality,
	$\displaystyle\biggl|\oint_{\partial D}h(z)\,dz\biggr|
	\le M\,\ell(\partial D).$
	On the other hand, since $f$ is analytic on the simply
	connected domain $D$, Cauchy's theorem gives
	$\oint_{\partial D}f(z)\,dz=0$, so
	\[
		\oint_{\partial D}h(z)\,dz
		=\oint_{\partial D}\bar z\,dz
		-\oint_{\partial D}f(z)\,dz
		=2iS(D).
	\]
	Hence
	\[
		2S(D)=|2iS(D)|\le M\,\ell(\partial D)
		\implies
		M\ge\frac{2S(D)}{\ell(\partial D)}.\qedhere
	\]
\end{solution}

\section{Exercise 3.63}

(Alternative proof of Corollary 3.15).
Prove Corollary 3.15 by Theorems 3.62 and 3.40.

\begin{solution}
	We show: on $\C^-:=\C\setminus(-\infty,0]$,
	$\Log z=\int_1^z\frac{d\zeta}{\zeta}$.

	\smallskip

	$\C^-$ is star-shaped with center $1$:
	for $z\in\C^-$, the segment
	$\{1+t(z-1):t\in[0,1]\}$ has imaginary part
	$t\,\Im z$ for $t>0$, so it never meets $(-\infty,0]$
	(when $\Im z=0$, $z>0$ and the segment stays on
	$(0,\infty)$).

	Since $1/\zeta$ is analytic on $\C^-$,
	Theorem~3.62 gives a primitive.
	By Theorem~3.40,
	$L(z):=\int_1^z d\zeta/\zeta$ is path-independent
	with $L'(z)=1/z$.

	Write $z=re^{i\theta}$, $r>0$, $\theta\in(-\pi,\pi)$.
	Integrate along $[1,r]\cup\{re^{it}:t\in[0,\theta]\}$:
	\[
		L(z)
		=\underbrace{\int_1^r\frac{dx}{x}}_{=\,\ln r}
		+\underbrace{\int_0^\theta
		\frac{ire^{it}}{re^{it}}\,dt}_{=\,i\theta}
		=\ln r+i\theta
		=\Log z.\qedhere
	\]
\end{solution}

\end{document}
